% Full candidate after original-source replay; not registered pending audits.
\subsection{A third independent iterated descent character}

\begin{statement}{C.31}{Open}{Candidate: three independent iterated parity characters}
There is an actual multiplier-one full gentle element $g_2$ such that
\[
                 \Delta(D(g_2))=0,\qquad
                 \Delta(D^2(g_2))=1.
 \tag{C.31.1}
\]
Consequently the continuous map
\[
 (\Delta,\Delta D,\Delta D^2):
       \widehat{\mathrm{GT}}_0\twoheadrightarrow C_2^3
 \tag{C.31.2}
\]
is onto, and is already onto on the multiplier-one subgroup.
These three characters are linearly independent and vanish
on the specified full Galois image. The arithmetic-normal quotient
therefore has an abelian quotient $C_2^3$.
This does not assert independence of all subsequent iterates, nor
does it determine the restriction of $\Delta$ to the pentagon group.
\end{statement}

\begin{proof}
We first solve one exact finite source-centralization problem, then
lift it to the full profinite modular group. Two normalizations are
essential: the first finite pure abelianization has exponent twelve
but its inertia has order twenty-four, so that normalization cannot
simply be omitted.

\paragraph{The actual marked finite source.}
Let $\Gamma=\widehat{C_2*C_3}=\langle s,r\mid s^2=r^3=1\rangle$,
$\tau=sr$, and $F=\ker(\Gamma\to S_3)$, with the old marking
$x=(rs)^2$, $y=\tau^2$. Write $G_1$ for the degree-$360$ group of
C.27, and $(\beta_1,A_1)$ for its actual corrected pure frames,
retained in \texttt{pbc\_hopf\_joint\_seed.json}. Put
\[
                  v_1=\beta_1,\qquad u_1=A_1^{-1}\beta_1.
 \tag{C.31.3}
\]
Retain the six evaluations of $q_1\varphi\sigma:F\to G_1$,
and also all six signed modulo-three and modulo-four evaluations
of $\varphi\sigma$, where $\sigma$ runs through the six geometric
peripheral permutations. Their actual Schreier-induced modular
representation has degree $2160+432=2592$. Denote its image by
$G=\Gamma/N$, its pure subgroup by $F_2$, and its generators by
$s,r$ again. The kernel is normal in $\Gamma$: the six substitutions
represent the geometric $S_3$ up to inner automorphisms, and the
induced representation retains exactly their simultaneous kernels.

In the ordered pairs
\[
           (s_1,\tau_1),(s_1,r_1^{-1}),
           (\tau_1,s_1),(\tau_1,r_1^{-1}),
           (r_1^{-1},s_1),(r_1^{-1},\tau_1),
\]
prescribe the six $G_1$ values
\[
 (v_1,u_1,v_1^{-1},u_1v_1^{-1},u_1^{-1},v_1u_1^{-1}),
 \tag{C.31.4}
\]
and prescribe identity in every first signed modulo-three and
modulo-four coordinate. The exact finite calculation places this
tuple in the derived subgroup of the actual joint word image.
Thus it is the value of an actual $F'$ word, not just a tuple in a
product. The word-image comparison with $F_2$ is an isomorphism;
let $f_2\in F_2'$ be the resulting element. The completed finite
calculations give
\[
\begin{array}{c|l}
 |G|&2^{146}3^{30}\\
 |F_2|&2^{145}3^{29}\\
 F_2^{\rm ab}&C_{12}^2\\
 \operatorname{ord}(x),\operatorname{ord}(y)&24,24\\
 G^{\rm ab}&C_6\\
 \alpha(s),\alpha(r)&s f_2^{-1},\ f_2r.
\end{array}
 \tag{C.31.5}
\]
Here $\alpha$ is an actual automorphism, fixing $\tau$ and admitting
pure $s$- and $r$-frames. The degree-$2592$ permutation generators,
all tuple coordinates, membership checks and marked comparison are
literal inputs of \texttt{pbc\_iterated\_joint.py} and its GAP
output. The finite portion of
\texttt{pbc\_hopf\_iterated\_joint\_sparse.g} completes the
permutation-to-polycyclic conversion and saves
\texttt{pbc\_iterated\_joint\_pc.json} before starting any Schur
calculation. The later memory failure of that old Schur attempt
does not enter this proof.

The completed triangular basis change
\texttt{pbc\_native\_adapt.py} retains every old-basis vector and
every new power and conjugation word. In the resulting $176$
generators $g_1,\ldots,g_{176}$, the marked elements are
\[
               r=g_1,\quad s=g_2,\quad
               B_s=g_7,\quad B_r=g_5.
 \tag{C.31.6}
\]
The fixed original source is
\texttt{pbc\_iterated\_shortcut\_v2\_original\_source.json},
with SHA256
\[
\texttt{95fb384593bfaf97d468bbc51f11cd797}
\quad\texttt{484a08cb3ec3b954961085b520ffc88}.
\]
It retains $71$ genuine pure $s$-centralizers and nine genuine pure
$r$-centralizers. A further finite-only permutation replay checks
all $176$ power and $15400$ conjugation clauses, before and after
the basis change, together with every marked frame and all these
pure centralizers, directly in the degree-$2592$ group. Its source
and output are \texttt{pbc\_iterated\_source\_permutation\_replay}.

\paragraph{The original integral central obstruction.}
Let $\mathcal F$ be the free discrete group on the $176$ marked
generators, $R=\ker(\mathcal F\to G)$, and
\[
              H=\mathcal F/[\mathcal F,R],\qquad
              K=R/[\mathcal F,R].
 \tag{C.31.7}
\]
The actual Hopf exact sequence identifies
$K\cap H'=H_2(G,\mathbb Z)$ with the torsion subgroup of the
finitely generated abelian group $K$; its remaining quotient is a
subgroup of the free abelian group $\mathcal F^{\rm ab}$.
All assertions obtained here as integral identities also hold on
profinite completion.

Use the canonical positive words $a,b$ for $\alpha(s),\alpha(r)$.
Each pair of words in the three differences
\[
       (B_s s)-(aB_s),\quad (B_r r)-(bB_r),\quad (ab)-(sr)
 \tag{C.31.8}
\]
has the same image in $G$. Interpreting each difference as the
central quotient of the two words and adding in $K$, their sum is
the actual product discrepancy
\[
             T=(B_s sB_s^{-1})(B_r rB_r^{-1})\tau^{-1}
                       \in K\cap H'.
 \tag{C.31.9}
\]
For a retained pure centralizer $c$ of $s$, put
$j_s(c)=csc^{-1}s^{-1}\in K\cap H'$. This is a homomorphism
on the finite centralizer: its values are central. It has exponent
dividing two, since $s^2$ is central in $H$. Likewise the
$r$-centralizer map $j_r$ has exponent dividing three.

C.42 proves that the three-primary part of the original $T$ is
zero. More precisely, its literal original-row identity in the
actual $26$-generator quotient is $T+3j_r(c_{74})=0$ after the
proved homology isomorphism, and $3j_r(c_{74})=0$ already in the
original centralization. We must solve the two-primary part with
the same original target.

\paragraph{A literal original-row identity in the $86$-generator quotient.}
Let $G_{86}$ be the actual quotient of C.43. C.42's integral
bound kills the original two-primary multiplier by $2^{14}$;
C.43's split surjection therefore gives the same bound on
$H_2(G_{86},\mathbb Z_2)$. For its $86$-generator free source,
the central relation module has $3741$ canonical power/conjugation
coordinates. If $K_{86}$ is its actual finitely generated abelian
central kernel, then
\[
 \operatorname{tors}(K_{86})_{(2)}\cap 2^{14}K_{86}=0.
 \tag{C.31.10}
\]
Indeed split $K_{86}$ abstractly into its free part and finite
torsion. Multiplication by $2^{14}$ kills the two-primary torsion,
and an element of the free part cannot be torsion. The odd-primary
part does not intersect the two-primary part.

The fresh calculation first recollects $9000$ consistency identities
using GAP's original positive-word source collector. The older
native-collected attempt failed its original-source comparison and
is not used. The successful fresh reducer gives an exact finite
DAG. Its complete expansion then replaces every auxiliary coordinate
or modulus leaf by its explicit Bezout combination of original
power/consistency rows. A separate literal GAP replay recollects all
$6446$ original leaves occurring in that expanded certificate.
The final identity is
\[
 T_{86}+\sum_i n_i j_s(c_i) +\sum_\ell m_\ell R_\ell
                      =2^{14}z\quad\text{in }\mathbb Z^{3741},
 \tag{C.31.11}
\]
where every $R_\ell$ is an actual source identity, every $c_i$ is
one of the original genuine pure centralizers, and $z$ is an
explicit integer vector. No independent power-word correction,
coordinate relation, modulus relation or $r$-frame correction
occurs in the final witness. The files
\texttt{pbc\_quotient86\_standard2\_literal\_original\_identity.json}
and the corresponding \texttt{\_replay.g/.out} retain every integer,
word and source comparison. The replay completed successfully on
all $6446$ leaves. The complete DAG and exact decoded-byte replay
are retained separately as reproducibility evidence.

Since every $j_s(c_i)$ has exponent two, only the parity of $n_i$
matters. In the original $176$-generator basis the nonzero corrections
are exactly
\[
\begin{split}
 \mathcal I=\{&6,16,27,35,39,41,43,49,\\
              &53,55,63,65,72,79,80,85\}.
\end{split}
 \tag{C.31.12}
\]
Here each $g_i$ for $i\in\mathcal I$ is one of the recorded pure
$s$-centralizers. The corrected quotient target is two-primary:
C.42 already kills the original target's three-primary part,
which remains zero after projection, and every added class is
two-torsion. Thus (C.31.10)--(C.31.11) prove that
\[
                   T+\sum_{i\in\mathcal I}j_s(g_i)
 \tag{C.31.13}
\]
maps to zero in $H_2(G_{86},\mathbb Z_2)$. This deduction uses
only genuine original identities; completeness of the searched
row set or of a guessed relation lattice is not assumed.

\paragraph{The two remaining original homology coordinates.}
C.43 identifies the kernel of
$H_2(G,\mathbb Z_2)\to H_2(G_{86},\mathbb Z_2)$ with
$(\bigwedge^2 A)_Q=C_2^2$, where
\[
       A=\ker(G\to Q)=C_2^{72},\qquad Q=G_{104},\qquad
       N_3=\ker(Q\to G_{86})=C_3^{18}.
\]
Here $A^{N_3}=A_{N_3}=0$, and the extension $G\to Q$ splits.
We need two actual central extensions of $G$ evaluating the two
coordinate forms; their mere existence would not determine the
value of (C.31.13).

The two invariant alternating forms $B_1,B_2$ of C.43 have ranks
$72$ and $36$. Their explicitly verified invariant quadratic
refinements are $q_1,q_2$. For either pair, take the bilinear form
$c$ with diagonal $q(e_i)$, upper triangular entries $B(e_i,e_j)$
and lower triangular entries zero. The finite group
\[
 E_q=A\times C_2,\qquad
 (v,z)(w,t)=(v+w,z+t+c(v,w))
 \tag{C.31.14}
\]
is a central extension of $A$. Its automorphisms fixing the
\emph{chosen} central $C_2$ kernel map onto $O(q)$ with kernel
$A^*$. For the rank-$36$ form the center is larger; no pointwise
fixing of that larger center is asserted.

Pull back this automorphism extension along the actual $Q$-action.
Odd-order averaging gives $(A^*)^{N_3}=0$ and
$H^2(Q,A^*)=0$, so it splits. Explicitly, for a matrix $M$ preserving
$q$, let $q_M$ have polar form $c(Mv,Mw)+c(v,w)$ and zero values on
the basis vectors. Its lifts are
\[
              (v,z)\longmapsto(Mv,z+q_M(v)+l(v)),\quad l\in A^*.
 \tag{C.31.15}
\]
The retained code first solves all order-three and commuting
relations of the $18$ actual $N_3$ generators, then uniquely lifts
the other quotient generators to normalize that complement.
There are $12312$ checked scalar equations of rank $1224$ in the
first step, and $1296$ equations of rank $72$ for each remaining
generator. All equations are derived from the original source
words and rechecked after solving them. Two automorphisms above
the same matrix differ by a linear form, so checking their
discrepancy on all basis vectors checks the whole identity.

The same normalizer construction gives an actual section of
$G\to Q$: a complement over $N_3$ has normalizer mapping onto
$Q$ by coprime complement conjugacy, and the kernel is
$A^{N_3}=0$. Every power and conjugation relation of this section
is additionally checked in the original finite $G$. Thus the
constructed actions give actual central covers $E_q\rtimes Q$ of
the actual $G$. This argument retains the $Q$ coordinate and does
not presume that its action on $A$ is faithful.

For clarity about the numerical convention, vectors are columns
and the right action of $g_i g_j$ is $M_jM_i$. Write a section lift
of $g_j$ as $g_j a_j$. On multiplying by the original $g_j$ in the
central cover, the coordinate update is
\[
 v'=M_jv+a_j,\qquad
 z'=z+q_{M_j}(v)+l_j(v)+c(M_jv,a_j).
 \tag{C.31.16}
\]
This evaluates all the positive words in (C.31.8) and every
centralizer commutator. A separate GAP calculation checks all
$83$ compared word pairs in the original finite group, not merely
their $A$-coordinates. The exact results are
\[
\begin{array}{c|c|c|c|c}
 \text{form}&\text{rank}&\text{three terms of }T&T&
             \text{all }s\text{- and }r\text{-corrections}\\\hline
 B_1&72&(0,0,0)&0&0\\
 B_2&36&(1,0,1)&0&0.
\end{array}
 \tag{C.31.17}
\]
The programs \texttt{pbc\_top\_two\_section},
\texttt{pbc\_top\_two\_cover\_actions}, and
\texttt{pbc\_top\_two\_cover\_values}, together with both original
GAP replays, retain the matrices, linear solutions, word values and
source hashes. Their independent reruns agree on every mathematical
field.

Restriction of a central extension of $A$ to $H_2(A,\mathbb Z)$
is its commutator pairing $B$. A change in quadratic diagonal lies
in the $\operatorname{Ext}^1$ part and does not affect this pairing.
The canonical inclusion in C.43 therefore identifies our two
cover evaluations on the kernel precisely with the two independent
forms $B_1,B_2$. An extension inflated from $Q$ evaluates zero on
that kernel by naturality. Consequently (C.31.17) detects every
remaining kernel element, and (C.31.13) is zero in the original
two-primary multiplier.

It is also zero at three: C.42 kills $T$ there, and each added
$s$-centralizer class is two-torsion. There are no other primary
parts, since $G$ has order a product of powers of two and three.
We have thus proved the single original integral identity
\[
                    T+\sum_{i\in\mathcal I}j_s(g_i)=0
                         \quad\text{in }K\cap H'.
 \tag{C.31.18}
\]
In particular the literal corrected frames are
\[
        \widetilde B_s=g_7\prod_{i\in\mathcal I}^{\rm increasing}g_i,
        \qquad \widetilde B_r=g_5.
 \tag{C.31.19}
\]
Their conjugates of $s,r$ have product exactly $sr$ in $H$.
The order of the displayed correction product is specified; the
centralizer commutator map makes its effect additive.

\paragraph{The modular source centralization and the full lift.}
Put $L=\langle g_2,g_1\rangle\subset H$. Its image is all of $G$,
so $H=LK$ and $L'=H'$ because $K$ is central. The images of
$g_2,g_1$ are independent basis vectors in $H^{\rm ab}=\mathbb Z^{176}$;
since they generate $L$, this proves $L^{\rm ab}=\mathbb Z^2$.
The actual free two-generator centralization maps onto $L$ and
induces the identity on the Hopf multiplier and an isomorphism on
free abelianizations. It is therefore an isomorphism. Hence
\[
              E=L/\langle g_2^2,g_1^3\rangle
                 =\Gamma/[N,\Gamma]
 \tag{C.31.20}
\]
after completion. The killed elements are central and have
independent nonzero abelian vectors, so their subgroup meets
$H'$ trivially. Since $G^{\rm ab}=C_6=\Gamma^{\rm ab}$,
the central kernel of $E\to G$ is exactly the finite Hopf
multiplier. This identifies the actual source centralization;
an arbitrary Schur cover is not being substituted for it.

Although the words in (C.31.19) were written in $H$, each can be
replaced by an element of $L$ with the same finite image. Their
difference is central and does not change any conjugation. Their
finite images are pure, so the chosen elements of $L$ are pure as
well. Equation (C.31.18) gives a homomorphism $\Gamma\to E$ that
conjugates $s,r$ by these frames and fixes $sr$ exactly. Its map on
the quotient is $\alpha$, so it sends $N$ into the central kernel
and kills $[N,\Gamma]$. It therefore factors through an endomorphism
of the actual $E$ in (C.31.20). Hopf naturality identifies its map on the
central kernel with $H_2(\alpha)$; it is consequently an
automorphism on kernel and quotient, hence on $E$.

Apply the modular central criterion proved by the finite trace and
refinement argument of P.3 and B.11. That proof uses only normality
of the retained source kernel, so it applies to this actual $N$.
It gives a full positive peripheral automorphism $\Phi$ of $\Gamma$
covering $\alpha$, retaining the two finite pure frames and the
identity cusp frame. An inner correction from $N$ makes
$\Phi(\tau)=\tau$ literally and changes none of the finite data.
The full torsion relations now give the full gentle two- and
three-cycle relations on restriction to $F$.

\paragraph{The first normalization and the common cusp factor.}
In the temporary marking $x_{\rm new}=\tau^{-2}$,
$y_{\rm new}=s x_{\rm new}s$, the conjugator
$s\Phi(s)$ has abelianization $(a,-a)$. Its finite value is
$f_2^{-1}\in F_2'$, so (C.31.5) implies
$a\in12\widehat{\mathbb Z}$. Normalize by the inner power
$x_{\rm new}^a$. Returning to the old marking replaces the finite
conjugator by
\[
                         f_2\longmapsto y^{-a}f_2x^a.
 \tag{C.31.21}
\]
This produces an actual normalized gentle pair $g_2=(1,f)$.
All first signed modulo-three and modulo-four values remain identity,
because all three corresponding peripheral matrices have orders
dividing twelve. The actual descended-frame formula of C.24
therefore gives $\Delta(D(g_2))=0$.

The change (C.31.21) need not be identity in $G$, whose inertia
has order twenty-four. Evaluate it in the next $G_1$ coordinates.
Since $s_1^a=r_1^a=1$, the first value and cycle-product value become
\[
       v'=\tau_1^{-a}\beta_1,\qquad
       u'=u_1=A_1^{-1}\beta_1.
 \tag{C.31.22}
\]
For the latter equality use the actual cycle $\rho$:
$\rho(y^{-a}f_2x^a)(y^{-a}f_2x^a)
=z^{-a}(\rho(f_2)f_2)x^a$, whose outer factors evaluate to one.
The retained actual pure frames of the first descended modular
automorphism are consequently
\[
              \beta_D=v'=\tau_1^{-a}\beta_1,\qquad
              A_D=v'(u')^{-1}=\tau_1^{-a}A_1.
 \tag{C.31.23}
\]
They have the \emph{same} left factor, a cusp-centralizing pure
power since $a$ is divisible by twelve. Remove it by a full inner
conjugation by $\tau^a$ before normalizing this descended class.
This does not change $D(g_2)$ as an outer class and restores exactly
the two C.27 finite frames, not just the same finite automorphism.
The certificate \texttt{pbc\_iterated\_normalization\_gate.py}
checks both possible values $a=0,12\pmod{24}$ against all six
actual permutation coordinates.

\paragraph{The second normalization and the signed character.}
With these restored frames the unnormalized pure conjugator of
$D(g_2)$ has the finite value $w$ of C.27 in the old marking.
It lies in $F_1'$, where $F_1$ is the pure subgroup of $G_1$.
Here $F_1^{\rm ab}=C_{12}^2$ and both inertia orders are twelve.
Thus the next abelian normalizing exponent is divisible by twelve,
and this second normalization is genuinely invisible in $G_1$.
The normalized actual pair representing $D(g_2)$ therefore retains
exactly C.27's six modulo-four residues and six projective
modulo-three identities.

Apply C.27's signed-three calculation to this actual full pair.
Its two allowed signed patterns have the same equianharmonic bit
one and Gaussian bit zero after the normalized intertwiner is
chosen. Equations (C.27.14)--(C.27.16), now applied to $D(g_2)$,
give $\Delta(D^2(g_2))=1$. This establishes (C.31.1) with both
normalizations and all retained pure frames accounted for.

Finally C.28 supplies the vector $(1,0,0)$, C.27 supplies a vector
$(b,1,c)$, and the present element supplies $(d,0,1)$. They span
$\mathbb F_2^3$, proving (C.31.2). Every character vanishes on
Galois because $D$ fixes the specified Galois image and $\Delta$
vanishes there. All stopped and failed calculations remain retained;
the positive assertion uses only the successful original-row,
actual-cover and full-lifting comparisons above.
\end{proof}
